\begin{textsample}{POS Dim 1 – 2013\_c5 – Score 10.00 – 2013\_c5/enade\_2013\_serv\_soc\_b0ca85ce74.txt}  \label{ex:f1_pos_082}
CAPÍTULO 5 DISTRIBUIÇÃO DOS CONCEITOS 5.1 \textbf{PANORAMA} NACIONAL DA DISTRIBUIÇÃO DOS CONCEITOS A Tabela 5.1 apresenta a quantidade e distribuição de cursos de Serviço Social participantes do ENADE/2013, por faixa de conceito e Grande Região.
A diferença entre os cursos \textbf{tabulados} neste capítulo e no capítulo 2 são os cursos sem conceito, em princípio, aqueles sem alunos concluintes que participassem da prova ou que tivessem tão somente um aluno concluinte14.
Observando -se os dados da Tabela 5.1, nota -se que, dos 282 cursos participantes, 139 ( 49,3 \% ) classificaram -se com conceito 3, o valor modal.
Este foi também o conceito modal em todas as regiões do Brasil ( entre 41,5 \% na região Sul e 67,2 \% na região Nordeste ), exceto na região Norte, onde a moda foi o conceito 2 ( 48,1 \% ).
O conceito 4 foi o segundo mais frequente em nível nacional ( 23,8 \%, correspondendo a 67 cursos ) e o conceito 2, o terceiro ( 18,8 \%, correspondendo a 53 cursos ).
Houve, ainda, 11 cursos ( 3,9 \% ) que receberam conceito 5 e outros seis cursos que receberam conceito 1.
Dos 282 cursos de Serviço Social, também seis cursos ficaram sem conceito ( SC ).
A região Norte participou com 27 cursos ou 9,6 \% do total nacional.
Nesta região, como já comentado, o conceito 2 foi atribuído a 13 cursos dos 27 participantes, o que \textbf{equivale} a 48,1 \% do total regional.
Os demais foram avaliados com conceito 3 ( nove cursos, correspondendo a 33,3 \% ), conceito 4 ( três cursos correspondendo a 11,1 \% ) e conceito 5 ( dois cursos correspondendo a 7,4 \% ).
Nenhum curso ficou sem conceito ( SC ) ou foi alocado ao conceito 1. 14 Estes cursos com somente um concluinte não foram considerados no capítulo 2.
A região Nordeste participou com 58 cursos ou 20,6 \% do total nacional.
Nesta região, 39 cursos, 67,2 \% em termos regionais, obtiveram conceito 3, o conceito modal para a região, como já comentado.
Os conceitos 2 e 4 foram atribuídos a oito cursos cada ( 13,8 \%, cada conceito ).
Já o conceito 5 foi atribuído a apenas, três cursos ( 5,2 \% ).
Nessa região, nenhum curso ficou sem conceito ( SC ) ou recebeu o conceito 1.
Como, também, já comentado, dos 129 cursos participantes da região Sudeste, 61 ( 47,3 \% ) obtiveram conceito 3, o conceito modal.
O conceito 1 foi atribuído a cinco cursos ( 3,9 \% ) e o conceito 2, a 18 ( 14,0 \% ).
O conceito 4 foi atribuído a 40 cursos ( 31,0 \% ) e dois outros ( 1,6 \% ) receberam o conceito 5.
Apenas três cursos ficaram sem conceito ( 2,3 \% ).
A região Sul também \textbf{contou} com cursos distribuídos em todas as faixas de conceitos.
A predominância do conceito 3 foi de 41,5 \%, como já comentado anteriormente, correspondentes a 22 dos 53 cursos participantes na região Sul.
O conceito 2 foi atribuído a dez cursos ( 18,9 \% ) e o conceito 4, a 13 cursos ( 24,5 \% ).
O conceito 1 foi recebido por um único curso ( 1,9 \% ) e o conceito 5, por quatro deles ( 7,5 \% ).
Apenas três dos cursos da região Sul ficaram sem conceito ( 5,7 \% ).
Como já comentado, também, de todos os 15 cursos participantes na região Centro-Oeste, oito receberam conceito 3, o conceito modal ( 53,3 \% ).
Nesta região nenhum curso ficou sem conceito ( SC ) ou recebeu os conceitos 1 ou 5.
Os demais cursos foram avaliados com conceito 2 ( quatro cursos, 26,7 \% ) e conceito 4 ( três cursos, 20,0 \% ).
O Gráfico 5.1 apresenta a distribuição do Conceito ENADE segundo Grande Região : \textbf{poligonais} mais à esquerda representam Grandes Regiões com pior distribuição, e \textbf{poligonais} mais à direita, Grandes Regiões com melhores conceitos.
Fica evidente que não houve um comportamento consistente entre as regiões, pois todas as \textbf{poligonais} relativas às diferentes regiões avaliadas Norte ( linha azul ), Nordeste ( linha verde ), Sudeste ( linha vermelha ) e Centro-Oeste ( linha preta ) se \textbf{entrecruzam}. 5.2 CONCEITOS POR CATEGORIA ADMINISTRATIVA E POR GRANDE REGIÃO A Tabela 5.2 apresenta a distribuição dos cursos participantes do ENADE/2013 de Serviço Social, por Categoria Administrativa, de acordo com os conceitos por eles alcançados, segundo as Grandes Regiões brasileiras.
Dos 282 cursos participantes, 57 ( 20,2 \% ) eram \textbf{ministrados} em Instituições \textbf{Públicas} e 225 ( 79,8 \% ), em Privadas.
De acordo com as informações da Tabela 5.2, em termos nacionais, das 11 IES que receberam o conceito 5, oito eram Públicas e três Privadas.
Dos 57 cursos participantes de IES Públicas, o conceito 4 foi o valor modal, atribuído a 19 cursos ( 33,3 \% ).
Entre os demais cursos de Instituições \textbf{Públicas} participantes, cinco obtiveram conceito 1 ( 8,8 \% da categoria ), oito receberam conceito 2 ( 14,0 \% ) e 16 foram alocados ao conceito 3 ( 28,1 \% ).
Como já comentado, os oito \textbf{restantes} receberam o conceito 5 ( 14,0 \% ).
Nesta categoria, apenas um dos cursos ficou sem conceito ( SC ).
Na rede privada, o conceito modal foi 3, com 123 cursos dos 225 da categoria.
Entre os demais cursos participantes, somente um recebeu conceito 1 e 45, conceito 2.
O conceito 4 foi atribuído a 48 cursos e o conceito 5, como já comentado, a três cursos.
Nesta Categoria Administrativa, apenas cinco dos cursos ficaram sem conceito ( SC ).
Na análise por região, observa -se que, na região Norte, as Instituições Privadas participaram com 20 cursos ( 74,1 \% do total regional, a menor proporção entre todas as regiões, juntamente com a região Nordeste ), dos quais nenhum ficou sem conceito ( SC ) ou recebeu os conceitos 1 ou 5.
O conceito modal para as Instituições Privadas na região foi 2, com 12 cursos, correspondendo a 60,0 \%.
O conceito 3 apresentou sete cursos, correspondendo a 35,0 \%.
Somente um curso foi alocado no conceito 4 ( 5,0 \% ).
As Instituições Públicas participaram com seis cursos alocados igualmente aos conceitos 3, 4 e 5 ( dois cursos, cada ), enquanto somente um curso recebeu conceito 2.
Nenhum curso ficou sem conceito ( SC ) ou recebeu conceito 1 nessa categoria na região.
Na região Nordeste, a rede privada concentrou 43 dos 58 cursos participantes, equivalentes a 74,1 \% do total da região, como já citado, apresentou juntamente com a região Norte, a menor proporção dentre todas as regiões.
Nesta Categoria Administrativa na região, predominaram os cursos com conceito 3, 35 cursos correspondendo a 81,4 \% dos cursos oferecidos por IES Privadas no Nordeste.
Apenas cinco cursos receberam o conceito 2 ( 11,6 \% ), enquanto o conceito 4 participou com três cursos ( 7,0 \% ).
Nenhum curso ficou sem conceito ( SC ) ou recebeu os conceitos 1 ou 5 nessa categoria e região.
As Instituições \textbf{Públicas} dessa região participaram com 15 cursos ( 25,9 \% ), dos quais cinco obtiveram conceitos 4, o valor modal.
Os demais foram avaliados com os conceitos 2 e 5 ( três cursos, cada ) e 3 ( quatro cursos ).
Nenhum curso ficou sem conceito ( SC ) ou recebeu conceito 1 nesta combinação de região e categoria.
Na região Sudeste, a proporção de cursos da rede privada, 85,3 \%, foi mais elevada do que nas demais regiões brasileiras, correspondendo a 110 dos 129 cursos participantes.
Nesta categoria, na região Sudeste, o conceito modal foi 3 ( 56 cursos ).
Nesta combinação de Categoria Administrativa e Grande Região, apenas três dos cursos não receberam conceito ( SC ).
Os demais foram avaliados com conceitos 1 e 5 ( um curso, cada ), conceito 2 ( 16 cursos ) e conceito 4 ( 33 cursos ).
Entre os 19 cursos em Instituições \textbf{Públicas} na região Sudeste, a categoria modal foi a que recebeu conceito 4 ( sete cursos ).
Os demais foram avaliados com os conceitos 1 ( quatro cursos ), 2 ( dois cursos ), 3 ( cinco cursos ) e 5 ( um curso ).
Na região Sudeste, nesta categoria, nenhum curso ficou sem conceito ( SC ).
As Instituições Privadas concentraram 40 dos 53 cursos participantes da região Sul, 75,5 \% do total regional.
Desses, 18 ficaram com conceito 3, o conceito modal.
Nesta combinação de Categoria Administrativa e Grande Região, dois cursos ficaram sem receber conceito ( SC ), enquanto nenhum recebeu conceito 1.
Os demais foram avaliados com conceito 2 ( oito cursos ), conceito 4 ( dez cursos ) e conceito 5 ( dois cursos ).
As Instituições \textbf{Públicas} na região Sul participaram com 13 cursos ( 24,5 \% ), dos quais somente um ficou sem conceito ( SC ).
O conceito modal também foi 3, com quatro cursos.
Além disso, os demais foram avaliados com conceito 1 ( um curso ), conceitos 2 e 5 ( dois cursos, cada ) e conceito 4 ( três cursos ).
Na região Centro-Oeste, 12 dos 15 cursos participantes eram de Instituições Privadas ( 80,0 \% em termos regionais ).
Destes, sete concentraram -se no conceito 3, conceito modal.
Os demais receberam conceito 2 ( quatro cursos ) e conceito 5 ( um curso ).
Nenhum curso ficou sem conceito ( SC ) ou recebeu conceitos 1 ou 5 nessa categoria.
Dos três cursos de Instituições \textbf{Públicas}, dois foram avaliados no conceito 4, conceito modal, e somente um curso foi alocado ao conceito 3.
Nesta região, nenhum curso de IES Pública ficou sem conceito ( SC ) ou recebeu os demais conceitos \textbf{restantes} nessa categoria e região. 5.3 CONCEITOS POR ORGANIZAÇÃO ACADÊMICA E POR GRANDE REGIÃO Na Tabela 5.3 encontra -se a distribuição dos conceitos atribuídos aos cursos participantes do ENADE/2013 na Área de Serviço Social, por Organização Acadêmica, segundo as Grandes Regiões brasileiras.
Dos 282 cursos participantes, 123 eram oferecidos em Faculdades, 122 em Universidades e os demais 37 em Centros Universitários.
Esta distribuição corresponde a, respectivamente, 43,6 \%, 43,3 \% e 13,1 \% dos cursos.
De acordo com os dados apresentados, de todos os 11 cursos avaliados com conceito 5, nove eram vinculados a Universidades.
Apenas dois cursos em Universidades ficaram sem conceito ( SC ).
Esse tipo de Organização Acadêmica teve o conceito 3 como modal, com 55 cursos ( 45,1 \% ).
Os demais cursos de Universidades avaliados, receberam os conceitos 1 ( quatro cursos ), 2 ( 15 cursos ), 4 ( 37 cursos ) e conceito 5 ( nove cursos, como já mencionado ).
Entre os cursos em Centros Universitários, o conceito modal também foi 3, com 16 cursos ( 43,2 \% ).
Neste tipo de Organização Acadêmica nenhum curso ficou sem conceito ( SC ).
Os outros cursos neste tipo de Organização Acadêmica receberam os conceitos 1 ( um curso ), 2 e 4 ( nove cursos, cada ) e conceito 5 ( dois cursos ).
Nas Faculdades, quatro dos 123 cursos ficaram sem conceito ( SC ) e 68 ( 55,3 \% ) receberam o mesmo conceito modal dos demais tipos de Organização Acadêmica, conceito 3.
Dos demais cursos neste tipo de Organização Acadêmica, um recebeu conceito 1, 29 conceito 2 e 21 conceito 4.
Nenhum curso recebeu o conceito 5 nessa categoria.
Considerando -se separadamente as regiões brasileiras, verifica -se que, na região Norte, as Universidades concentraram dez dos 27 cursos participantes.
Dos cursos em Universidades, nenhum curso ficou sem conceito ( SC ) ou recebeu conceito 1.
Quatro cursos ficaram alocados igualmente aos conceitos 4 e 5 ( dois cursos, cada ) e cinco receberam o conceito modal, 3.
Os Centros Universitários da região Norte foram representados por quatro cursos, os quatro distribuídos igualmente aos conceitos 2 e 3.
Nenhum curso ficou sem conceito ( SC ) ou recebeu os demais conceitos \textbf{restantes} nessa categoria e região.
As Faculdades participaram com 13 cursos na região Norte.
Nenhum ficou sem conceito ( SC ) ou recebeu conceitos 1 ou 5.
Os demais receberam conceito 2 ( dez cursos, conceito modal ), conceito 3 ( dois cursos ) e conceito 4 ( um curso ).
Na região Nordeste, as Universidades participaram com 23 dos 58 cursos na Área de Serviço Social da região.
Nenhum dos cursos em Universidades no Nordeste ficou sem conceito ( SC ) ou recebeu conceito 1.
O conceito modal foi 3, com 11 cursos.
Os demais receberam conceito 2 ( quatro cursos ), 4 ( cinco cursos ) e conceito 5 ( três cursos ).
Os Centros Universitários \textbf{contaram} com dois cursos participantes na região Nordeste, todos alocados igualmente aos conceitos 3 e 4 ( um curso, cada ).
Nenhum curso ficou sem conceito ( SC ) ou recebeu os demais conceitos \textbf{restantes} nessa categoria e região.
As Faculdades foram representadas por 33 cursos na região Nordeste, e o conceito modal foi 3, com 27 cursos.
Dos cursos \textbf{restantes}, quatro receberam conceito 2 e dois, conceito 4.
Nessa categoria e região, nenhum curso ficou sem conceito ( SC ) ou recebeu os conceitos 1 ou 5.
Na região Sudeste, as Universidades concentraram 50 dos 129 cursos da região.
Entre os cursos em Universidades na região, os conceitos modais foram 3 e 4 com 20 cursos, cada ( \textbf{bimodal} ).
Nenhum curso ficou sem conceito ( SC ), enquanto quatro cursos receberam conceito 1.
Os demais cursos receberam os conceitos 2 ( cinco cursos ) e 5 ( um curso ).
Os Centros Universitários participaram com 26 cursos na região Sudeste, dos quais 12 obtiveram conceito modal, 3, e nenhum ficou sem conceito ( SC ).
Os demais receberam os conceitos 1 e 5 ( um curso, cada ), 2 ( quatro cursos ) e 4 ( oito cursos ).
As Faculdades foram representadas por 53 cursos na região Sudeste, que se distribuíram nos conceitos 2 ( nove cursos ), 3 ( 29 cursos, conceito modal ) e 4 ( 12 cursos ).
Três cursos ficaram sem conceito ( SC ), enquanto nenhum recebeu os conceitos 1 ou 5.
Dos 53 cursos da região Sul, 31 eram de Universidades, para os quais o conceito modal foi 3, com 14 cursos.
Nesse tipo de organização, nenhum dos cursos recebeu o conceito 1.
Os demais receberam os conceitos 2 ( quatro cursos ), 4 ( oito cursos ) e 5 ( três cursos ).
Apenas dois cursos ficaram sem conceito ( SC ) nessa categoria e região.
Os Centros Universitários da região Sul tiveram dois dos três cursos participantes no conceito modal, 2, enquanto, somente um curso recebeu conceito 5.
Nenhum curso ficou sem conceito ( SC ) ou recebeu os demais conceitos \textbf{restantes} nessa categoria.
Foram 19 os cursos vinculados a Faculdades na região Sul, e somente um destes ficou sem conceito ( SC ).
De todos os cursos, oito receberam conceito 3, o conceito modal.
Os demais receberam os conceitos 1 ( um curso ), 2 ( quatro cursos ) e 4 ( cinco cursos ).
Nessa categoria e região, nenhum curso obteve o conceito 5.
Na região Centro-Oeste, oito dos 15 cursos eram de Universidades.
Nesse tipo de organização, nenhum curso ficou sem conceito ( SC ) ou recebeu os conceitos 1 ou 5.
O conceito modal foi 3, com cinco cursos.
Os outros cursos obtiveram os conceitos 2 ( um curso ) e 4 ( dois cursos ).
Os Centros Universitários da região Centro-Oeste \textbf{contaram} com dois cursos, todos foram alocados igualmente aos conceitos 2 e 3 ( um curso, cada ).
Nenhum curso ficou sem conceito ( SC ) ou recebeu os demais conceitos \textbf{restantes} nessa categoria.
Dos cinco cursos em Faculdades na região Centro-Oeste, nenhum ficou sem conceito ( SC ) ou recebeu conceitos 1 ou 5.
Quatro cursos foram distribuídos igualmente, também, nos conceitos 2 e 3 ( dois cursos, cada ), enquanto somente um curso recebeu conceito nessa categoria e região.
O Gráfico 5.3 apresenta a distribuição do Conceito ENADE segundo a Organização Acadêmica da IES.
Os cursos em Universidades ( linha azul ) apresentam uma \textbf{poligonal} mais à direita do que as demais, \textbf{denotando} uma melhor distribuição de Conceitos ENADE.
No outro extremo temos os cursos em Faculdades ( linha vermelha ) com a \textbf{poligonal} mais à esquerda e consequentemente uma distribuição de Conceitos ENADE com menores valores.
Cursos em Centros Universitários ( linha verde ) apresentam uma situação intermediária.

% matched lemmas: bimodal, contar, denotar, entrecruzar, equivaler, ministrar, panorama, poligonal, pública, restante, tabular
\end{textsample}
